\documentclass[12pt,a4paper]{article}

\usepackage[utf8]{inputenc}
\usepackage[T1]{fontenc}
\usepackage[brazil]{babel}
\usepackage{amsmath}
\usepackage{geometry}

\usepackage{indentfirst}
\setlength{\parindent}{0.8cm}

\usepackage[
    backend=biber,
    style=abnt
]{biblatex}

\addbibresource{references.bib}

\geometry{
    top=2.5cm,
    bottom=2.5cm,
    left=2.5cm,
    right=2.5cm
}

\title{\textbf{Otimização de Contratações Contextuais em Elencos de Futebol}}
\author{Lucio Vargas de Albuquerque Nunes}
\date{GCC118 -- Programação Matemática \\ 2026/2}

\begin{document}

\maketitle

\section*{Introdução}

\section{Introdução}

\section{Introdução}

A composição de elencos no futebol profissional constitui um problema de tomada de decisão que envolve a seleção de jogadores sob restrições financeiras, esportivas e estruturais \parencite{pantuso2017}. A escolha de novos atletas não depende exclusivamente de suas características individuais, mas também de sua compatibilidade com os demais integrantes da equipe e de sua contribuição potencial para o desempenho coletivo \parencite{jarvandi2013}. Nesse contexto, técnicas de Programação Matemática têm sido empregadas para representar o problema de composição de elencos, permitindo identificar combinações de jogadores que maximizem determinados critérios de desempenho, respeitando restrições orçamentárias, posicionais e regulatórias \parencite{pantuso2017,dantas2025}.

A integração entre modelos preditivos e técnicas de otimização possibilita avaliar contratações não apenas por indicadores individuais, mas também por seus efeitos estimados sobre o desempenho de uma equipe. Nessa perspectiva, \textcite{nunes2026} desenvolveu uma metodologia que utiliza estatísticas individuais de jogadores para prever a pontuação de clubes ao longo de uma temporada, empregando posteriormente essas previsões como função objetivo de algoritmos de otimização. Como continuidade dessa abordagem, o presente trabalho propõe investigar uma formulação de Programação Linear Inteira (PLI) baseada em ganhos previamente estimados para possíveis substituições de jogadores, permitindo a aplicação e comparação de diferentes métodos de resolução.

A base de dados utilizada é proveniente do FBref, plataforma que disponibiliza estatísticas históricas de jogadores e equipes de futebol \parencite{fbref}. Os dados foram coletados e estruturados por \textcite{nunes2026}, abrangendo cinco temporadas, de 2019--2020 a 2023--2024, da Premier League inglesa e da Serie A italiana. O conjunto contempla indicadores de participação, características individuais e estatísticas de desempenho ofensivo e defensivo dos atletas. Além disso, a metodologia original incorpora uma função sintética de valor de mercado baseada na idade e em indicadores estatísticos dos jogadores, permitindo representar restrições orçamentárias mesmo na ausência de valores reais de transferência para todos os atletas \parencite{nunes2026}. A consideração de características individuais e de desempenho na estimativa de custos encontra respaldo em estudos sobre os determinantes econômicos das transferências no futebol profissional \parencite{dobson1999}.

Entretanto, a aplicação de modelos de otimização depende não apenas da disponibilidade dos dados, mas também de sua qualidade, consistência e adequação às variáveis e restrições do problema. A necessidade de organização, integração e pré-processamento das estatísticas individuais dos jogadores é evidenciada na metodologia desenvolvida por \textcite{nunes2026}. No contexto da composição de elencos, aspectos como a quantidade de jogadores elegíveis, a disponibilidade de atributos estatísticos, a definição de candidatos e vagas e a distribuição dos custos de contratação influenciam a construção das instâncias de otimização e a definição de suas restrições \parencite{pantuso2017,dantas2025}. Dessa forma, torna-se necessário caracterizar a base de dados e examinar suas condições de utilização antes da implementação da formulação matemática.

Assim, esta etapa do trabalho tem como objetivo realizar uma análise exploratória dos dados previamente coletados e estruturados por \textcite{nunes2026}, contemplando a caracterização da base, a avaliação de sua qualidade e o dimensionamento das instâncias de otimização. Busca-se identificar possíveis limitações, examinar a disponibilidade das informações necessárias e avaliar sua adequação à construção de cenários de contratação sob restrições orçamentárias e estruturais. Os resultados dessa análise servirão de fundamento para as etapas posteriores de formulação matemática, desenvolvimento de métodos heurísticos e comparação experimental das abordagens de resolução.

\section*{Caracterização da Base de Dados}

\section*{Pré-processamento e Qualidade dos Dados}

\section*{Construção das Instâncias de Otimização}

\section*{Análise de Viabilidade}

\section*{Conclusão}

\end{document}